% Bài giữa kỳ (cá nhân). Layout shared with the other reports of the course in
% ../../shared/. Figures come from ../figures/, written by ../make_figures.py.
% Compile with XeLaTeX: latexmk -xelatex report.tex
\documentclass[12pt,a4paper]{report}

\newcommand{\sharedir}{../../shared}
\newcommand{\figdir}{../figures}
\input{\sharedir/layout.tex}

% ============ TITLE PAGE FIELDS ============
\newcommand{\reporttitle}{BÀI TẬP GIỮA KỲ}
\newcommand{\reportname}{Ước lượng (học) mô hình hồi quy tuyến tính Bayes \\
  và sự phụ thuộc vào đặc điểm của mẫu dữ liệu học}
\newcommand{\studentlabel}{Học viên thực hiện}
\newcommand{\studentname}{Bùi Quang Huy (25007877)}

\begin{document}

\makereporttitlepage

\chapter*{Tóm tắt}
\addcontentsline{toc}{chapter}{Tóm tắt}

Ước lượng mô hình hồi quy tuyến tính Bayes là xác định phân phối hậu nghiệm
của các hệ số và phương sai nhiễu từ dữ liệu và phân phối tiên nghiệm. Bài viết
trình bày quy trình gồm bảy bước: chuẩn bị dữ liệu, đặc tả hàm hợp lý, chọn tiên
nghiệm, tính hậu nghiệm, chẩn đoán hội tụ, kiểm tra mô hình, suy luận và dự đoán.
Trọng tâm là phân tích mối liên hệ giữa đặc điểm của mẫu và quá trình ước lượng.
Trong mô hình sai số chuẩn đang xét, dữ liệu ảnh hưởng đến hậu nghiệm thông qua
$n$, $X\trans X$, $X\trans y$ và $y\trans y$. Các công thức và mô phỏng được
sử dụng để minh hoạ vai trò của cỡ mẫu, mức nhiễu, phạm vi và mức độ phân tán
của biến giải thích, đa cộng tuyến, thang đo và ngoại lai. Phần ứng dụng đối
chiếu quy trình với kết quả phân tích 93 mẫu xe năm 1993 trong bài cuối kỳ
của nhóm.

\tableofcontents
\pagebreak

\listoffigures
\listoftables
\newpage

% ===========================================================================
\chapter{Bài toán học mô hình hồi quy tuyến tính Bayes}
\label{sec:problem}

\section{Mô hình}

Cho mẫu học $\mathcal{D} = \{(x_i, y_i)\}_{i=1}^{n}$, với $x_i \in \R^{p}$ là
vector biến giải thích (phần tử đầu bằng 1 ứng với hệ số chặn) và $y_i \in \R$
là biến phản hồi. Ghép thành ma trận thiết kế $X \in \R^{n \times p}$ và vector
$y \in \R^{n}$, mô hình hồi quy tuyến tính chuẩn là
%
\begin{equation}
  \label{eq:model}
  y \mid X, \beta, \sigma^2 \sim \N\bigl(X\beta,\; \sigma^2 I_n\bigr).
\end{equation}
%
Cách tiếp cận Bayes coi $\theta = (\beta, \sigma^2)$ là biến ngẫu nhiên với phân
phối tiên nghiệm $p(\theta)$. \emph{Học} mô hình nghĩa là tính phân phối hậu
nghiệm
%
\begin{equation}
  \label{eq:posterior}
  p(\beta, \sigma^2 \mid y, X) =
  \frac{p(y \mid X, \beta, \sigma^2)\, p(\beta, \sigma^2)}
       {\int p(y \mid X, \beta, \sigma^2)\, p(\beta, \sigma^2)\, d\beta\, d\sigma^2}.
\end{equation}
%
Kết quả hồi quy thông thường được tóm tắt bằng ước lượng điểm $\hat\beta$ cùng
sai số chuẩn và khoảng tin cậy. Kết quả ước lượng Bayes được biểu diễn bằng
phân phối hậu nghiệm trên không gian tham số. Mẫu số của
\eqref{eq:posterior}, gọi là bằng chứng (evidence) hay hàm hợp lý biên, là một
tích phân nhiều chiều; việc tính tích phân này là một khó khăn chính trong suy luận Bayes \cite{gelman2013bda}.

\section{Hàm hợp lý và các thống kê đủ}
\label{sec:sufficient}

Lấy logarit hàm hợp lý của \eqref{eq:model} và khai triển,
%
\begin{equation}
  \label{eq:loglik}
  \log p(y \mid X, \beta, \sigma^2)
  = -\frac{n}{2}\log(2\pi\sigma^2)
    - \frac{1}{2\sigma^2}\Bigl( y\trans y - 2\beta\trans X\trans y
    + \beta\trans X\trans X \beta \Bigr).
\end{equation}
%
Hàm hợp lý chỉ phụ thuộc vào dữ liệu qua bốn đại lượng: $n$, $X\trans X$,
$X\trans y$ và $y\trans y$. Đó là các \emph{thống kê đủ} của mô hình. Trong mô hình \eqref{eq:model}, với tiên nghiệm cố định, ảnh hưởng của dữ liệu
đến hậu nghiệm được thể hiện thông qua bốn đại lượng này.
Chương~\ref{sec:dependence} lần lượt chỉ ra mỗi đặc điểm của mẫu làm thay đổi đại
lượng nào.

% ===========================================================================
\chapter{Các bước ước lượng mô hình}
\label{sec:steps}

Quá trình học được chia thành bảy bước. Thứ tự này cũng là quy trình nhóm dùng
trong bài cuối kỳ.

\section{Bước 1: Chuẩn bị dữ liệu}
\label{sec:step-data}

Làm sạch (giá trị thiếu, lỗi nhập liệu, đơn vị), chọn biến, biến đổi để quan hệ
gần tuyến tính và sai số gần chuẩn (chẳng hạn lấy logarit biến phản hồi lệch
phải), và chuẩn hoá các biến giải thích về trung bình 0, độ lệch chuẩn 1. Chuẩn
hoá không làm thay đổi mô hình, nhưng làm cho phân phối tiên nghiệm có cùng ý nghĩa với mọi hệ
số và giúp bước 4 hội tụ nhanh hơn, như mục~\ref{sec:dep-scale} giải thích.

\section{Bước 2: Đặc tả hàm hợp lý}

Chọn phân phối của sai số. Mặc định là chuẩn như \eqref{eq:model}; nếu dữ liệu có
ngoại lai thì thay bằng Student-t,
$y_i \sim t_\nu(x_i\trans\beta, \sigma)$, với $\nu$ được ước lượng từ dữ liệu. Phương sai
thay đổi theo nhóm, sai số tương quan theo thời gian hay cấu trúc phân cấp cũng
được đưa vào ở bước này.

\section{Bước 3: Chọn phân phối tiên nghiệm}
\label{sec:step-prior}

Ba lựa chọn thường dùng được trình bày dưới đây.

\begin{itemize}
  \item \textbf{Phân phối tiên nghiệm liên hợp} Normal--Inverse-Gamma (NIG):
    $\beta \mid \sigma^2 \sim \N(m_0, \sigma^2 V_0)$,
    $\sigma^2 \sim \text{IG}(a_0, b_0)$. Hậu nghiệm cùng họ, có công thức đóng.
  \item \textbf{Phân phối tiên nghiệm bán liên hợp}: $\beta \sim \N(m_0, V_0)$ độc lập với
    $\sigma^2 \sim \text{IG}(a_0, b_0)$. Không có dạng đóng cho hậu nghiệm đồng
    thời, nhưng các phân phối điều kiện có dạng đóng, nên dùng được Gibbs.
  \item \textbf{Phân phối tiên nghiệm tuỳ ý} như $\sigma \sim \text{Half-Normal}(1)$ hay
    $\beta_j \sim t_3(0, 2{,}5)$, cần MCMC tổng quát như NUTS.
\end{itemize}

Về mức thông tin, phân phối tiên nghiệm có thể phẳng, yếu thông tin hoặc tập trung.
\emph{Kiểm tra dự đoán tiên nghiệm (prior predictive check)} được thực hiện bằng
cách sinh dữ liệu mô phỏng từ phân phối tiên nghiệm và mô hình,
sau đó đánh giá mức độ hợp lý của dữ liệu mô phỏng trước khi sử dụng dữ liệu
quan sát để cập nhật hậu nghiệm.

\section{Bước 4: Tính phân phối hậu nghiệm}
\label{sec:step-posterior}

\subsection{Trường hợp $\sigma^2$ đã biết}

Với phân phối tiên nghiệm $\beta \sim \N(m_0, V_0)$, tích của hàm hợp lý \eqref{eq:loglik} và
phân phối tiên nghiệm là hàm mũ của một dạng toàn phương theo $\beta$. Hoàn thành bình phương
(phụ lục~\ref{sec:derivation}) cho hậu nghiệm chuẩn $\beta \mid y \sim
\N(m_n, V_n)$ với
%
\begin{equation}
  \label{eq:known-sigma}
  V_n = \Bigl(V_0^{-1} + \tfrac{1}{\sigma^2} X\trans X\Bigr)^{-1},
  \qquad
  m_n = V_n \Bigl(V_0^{-1} m_0 + \tfrac{1}{\sigma^2} X\trans y\Bigr).
\end{equation}
%
Ma trận độ chính xác hậu nghiệm $V_n^{-1}$ là tổng của ma trận độ chính xác
tiên nghiệm $V_0^{-1}$ và thành phần thông tin do dữ liệu cung cấp
$X\trans X / \sigma^2$; trung bình hậu nghiệm là trung bình
có trọng số của $m_0$ và ước lượng OLS $\hat\beta = (X\trans X)^{-1}X\trans y$,
vì $X\trans y = X\trans X\, \hat\beta$:
%
\begin{equation}
  \label{eq:weighted}
  m_n = V_n \Bigl( V_0^{-1} m_0 + \tfrac{1}{\sigma^2} X\trans X\, \hat\beta \Bigr).
\end{equation}

\subsection{Trường hợp liên hợp NIG}

Khi $\sigma^2$ cũng chưa biết và phân phối tiên nghiệm là NIG, hậu nghiệm vẫn là NIG với
%
\begin{equation}
  \label{eq:nig}
  \begin{aligned}
    V_n &= \bigl(V_0^{-1} + X\trans X\bigr)^{-1}, &
    m_n &= V_n \bigl(V_0^{-1} m_0 + X\trans y\bigr), \\
    a_n &= a_0 + \tfrac{n}{2}, &
    b_n &= b_0 + \tfrac{1}{2}\bigl(y\trans y + m_0\trans V_0^{-1} m_0
           - m_n\trans V_n^{-1} m_n\bigr).
  \end{aligned}
\end{equation}
%
Lấy tích phân theo $\sigma^2$, phân phối biên của $\beta$ là phân phối $t$ đa
biến với $2a_n$ bậc tự do, tâm $m_n$ và ma trận tỉ lệ $(b_n/a_n) V_n$
\cite{hoff2009first}. Số hạng $b_n$ chứa tổng bình phương phần dư: dữ liệu càng
nhiễu thì $b_n$ càng lớn, hậu nghiệm của $\sigma^2$ dịch về phía các giá trị lớn hơn và hậu nghiệm của
$\beta$ rộng ra.

\subsection{Lấy mẫu Gibbs với phân phối tiên nghiệm bán liên hợp}

Với phân phối tiên nghiệm bán liên hợp, hậu nghiệm đồng thời không có dạng đóng, nhưng hai phân
phối điều kiện đầy đủ thì có:
%
\begin{equation}
  \label{eq:gibbs}
  \begin{aligned}
    \beta \mid \sigma^2, y &\sim \N(m_n, V_n)
      \quad \text{với } m_n, V_n \text{ như } \eqref{eq:known-sigma}, \\
    \sigma^2 \mid \beta, y &\sim \text{IG}\Bigl(a_0 + \tfrac{n}{2},\;
      b_0 + \tfrac{1}{2}\norm{y - X\beta}_2^2\Bigr).
  \end{aligned}
\end{equation}
%
Thuật toán Gibbs \cite{geman1984gibbs} lần lượt lấy mẫu $\beta$ và $\sigma^2$ từ
hai phân phối điều kiện này. Sau giai đoạn khởi động, các mẫu thu được dùng
để xấp xỉ phân phối hậu nghiệm.

\subsection{MCMC tổng quát: Metropolis--Hastings và NUTS}

Khi khó tính mẫu số của \eqref{eq:posterior}, có thể sử dụng mật độ hậu nghiệm
chưa chuẩn hoá $\tilde p(\theta) = p(y \mid \theta)\, p(\theta)$. Mật độ này
được xác định đến một hằng số nhân. Thuật toán Metropolis--Hastings
\cite{hastings1970} đề xuất $\theta' \sim q(\theta' \mid \theta)$ rồi chấp nhận
với xác suất
%
\begin{equation}
  \label{eq:mh}
  \alpha = \min\Bigl(1,\;
    \frac{\tilde p(\theta')\, q(\theta \mid \theta')}{\tilde p(\theta)\, q(\theta' \mid \theta)}\Bigr),
\end{equation}
%
trong đó hằng số chưa biết triệt tiêu ở tỉ số. Bước đề xuất ngẫu nhiên khiến
chuỗi di chuyển chậm và tự tương quan cao khi tham số tương quan với nhau. NUTS
\cite{hoffman2014nuts}, một biến thể của Hamiltonian Monte Carlo, dùng gradient
của $\log \tilde p$ để xác định hướng và độ dài bước di chuyển phù hợp với hình dạng của hậu nghiệm, và là
thuật toán mặc định trong Stan và PyMC.

\subsection{Xấp xỉ}

Khi dữ liệu rất lớn, có thể xấp xỉ hậu nghiệm bằng phân phối chuẩn quanh điểm cực
đại (xấp xỉ Laplace) hoặc bằng suy diễn biến phân, nhằm giảm thời gian tính toán, đồng thời có thể làm giảm độ chính xác của xấp xỉ
\cite{bishop2006prml}.

\section{Bước 5: Chẩn đoán hội tụ}

Cần đánh giá hội tụ trước khi sử dụng kết quả MCMC để suy luận. Các tiêu chí thường
dùng là $\Rhat$ dựa trên phép chuẩn hoá thứ hạng dưới \num{1.01} \cite{vehtari2021rhat}, cỡ mẫu
hiệu dụng (ESS) bulk và tail đủ lớn (thường trên 400), không phát hiện hiện tượng phân kỳ khi lấy mẫu (divergence) với NUTS.
Cần xem thêm biểu đồ vết và biểu đồ thứ hạng để đánh giá tính ổn định và mức độ
nhất quán giữa các chuỗi. Với lời giải dạng đóng thì
không cần bước này, và đó là một lý do để đối chiếu MCMC với dạng đóng khi có thể.

\section{Bước 6: Kiểm tra và so sánh mô hình}

Kiểm tra dự đoán hậu nghiệm (posterior predictive check) sử dụng dữ liệu mô phỏng từ phân phối dự đoán hậu nghiệm
để so sánh với dữ liệu quan sát, theo mật độ hay theo các thống kê như giá trị lớn nhất, nhỏ nhất. So
sánh các mô hình ứng viên dùng ELPD ước lượng qua PSIS-LOO hoặc WAIC
\cite{vehtari2017loo}; tham số Pareto $k$ của PSIS-LOO còn chỉ ra các quan sát
có ảnh hưởng lớn tới hậu nghiệm.

\section{Bước 7: Suy diễn và dự đoán}

Từ hậu nghiệm, mọi đại lượng quan tâm được tính trực tiếp: trung bình, khoảng tin
cậy hậu nghiệm (HDI), xác suất $P(\beta_j > 0 \mid y)$. Dự đoán tại điểm mới
$x^\ast$ dùng phân phối dự đoán hậu nghiệm; với $\sigma^2$ đã biết,
%
\begin{equation}
  \label{eq:predictive}
  y^\ast \mid x^\ast, y \sim
  \N\bigl(x^{\ast\top} m_n,\; \sigma^2 + x^{\ast\top} V_n\, x^\ast\bigr).
\end{equation}

% ===========================================================================
\chapter{Quá trình học phụ thuộc vào đặc điểm nào của mẫu}
\label{sec:dependence}

Trong mô hình sai số chuẩn ở mục~\ref{sec:sufficient}, với tiên nghiệm cố định,
dữ liệu ảnh hưởng đến hậu nghiệm thông qua $n$, $X\trans X$, $X\trans y$ và
$y\trans y$. Các đại lượng này phản ánh cỡ mẫu, mức độ phân tán và tương quan
của các biến giải thích, cũng như mối liên hệ của chúng với biến phản hồi.
Các mục dưới đây phân tích vai trò của từng đặc điểm đối với quá trình ước
lượng và dự đoán. Trừ khi có ghi chú khác, các mô phỏng sử dụng mô hình sinh
dữ liệu $y = 1 + 2x + \varepsilon$ và lời giải dạng đóng
\eqref{eq:known-sigma}; mã nguồn tại \texttt{bai-giua-ky/make\_figures.py}.

\section{Cỡ mẫu}
\label{sec:dep-n}

$X\trans X$ là tổng của $n$ ma trận $x_i x_i\trans$, nên lớn lên tỉ lệ với $n$.
Theo \eqref{eq:known-sigma}, phương sai hậu nghiệm giảm theo $1/n$, độ lệch chuẩn
theo $1/\sqrt{n}$, còn phần đóng góp của phân phối tiên nghiệm $V_0^{-1}$ không đổi và nhỏ dần
tương đối. Hình~\ref{fig:sequential} cho thấy hậu nghiệm của hệ số góc từ một
phân phối tiên nghiệm rộng thu hẹp dần về giá trị thật khi $n$ tăng từ 2 lên 250. Khi $n$ rất
lớn, hậu nghiệm xấp xỉ $\N\bigl(\hat\beta, \sigma^2 (X\trans X)^{-1}\bigr)$, tức
trùng với kết quả của hồi quy thông thường; đây là nội dung của định lý
Bernstein--von Mises \cite{gelman2013bda}.

\resultfig[0.95]{fig_sequential}{Cập nhật tuần tự: hậu nghiệm của hệ số góc $\beta_1$
  khi cỡ mẫu tăng (trái), và trung bình hậu nghiệm cùng khoảng 95\% theo $n$ trên
  thang logarit (phải). Phân phối tiên nghiệm $\N(0, 2^2)$, $\sigma = 1$.}{fig:sequential}

Cỡ mẫu còn quyết định cách tổ chức việc học. Vì hậu nghiệm sau $n$ quan sát là
phân phối tiên nghiệm cho quan sát thứ $n+1$, có thể học tuần tự: cập nhật $V_n$ và $m_n$ bằng
\eqref{eq:known-sigma} mỗi khi có lô dữ liệu mới, và với cùng mô hình và tiên nghiệm ban đầu, kết quả tương đương với
ước lượng trên toàn bộ dữ liệu.

Ở bước 3, cỡ mẫu ảnh hưởng đến vai trò tương đối của tiên nghiệm. Với một biến giải thích đã
chuẩn hoá và phân phối tiên nghiệm $\N(0, \tau^2)$, trung bình hậu nghiệm của hệ số góc là
$w\,\hat\beta$ với
%
\begin{equation}
  \label{eq:shrinkage}
  w = \frac{n}{n + \sigma^2/\tau^2}.
\end{equation}
%
Tỉ số $\sigma^2/\tau^2$ đóng vai trò ``số quan sát tương đương'' của phân phối tiên nghiệm.
Hình~\ref{fig:prior-data} cho thấy với tiên nghiệm tập trung $\tau = 0{,}05$ (tương đương
400 quan sát), sau \num{1000} điểm dữ liệu trung bình hậu nghiệm vẫn chỉ đạt
khoảng 70\% giá trị thật; còn với phân phối tiên nghiệm gần phẳng $\tau = 100$, trung bình hậu
nghiệm dao động mạnh ở vài điểm đầu, vì khi đó gần như chỉ có dữ liệu quyết định.
Mẫu càng nhỏ thì bước chọn phân phối tiên nghiệm và phân tích độ nhạy càng quan trọng.

\resultfig[0.95]{fig_prior_vs_data}{Ảnh hưởng của tiên nghiệm và cỡ mẫu đến hậu nghiệm: trung bình hậu nghiệm
  của $\beta_1$ theo cỡ mẫu với bốn mức độ tập trung của tiên nghiệm $\N(0, \tau^2)$ (trái),
  và trọng số của dữ liệu $w$ theo \eqref{eq:shrinkage} (phải).}{fig:prior-data}

\section{Mức nhiễu của dữ liệu}
\label{sec:dep-noise}

Trong \eqref{eq:known-sigma}, dữ liệu đi vào dưới dạng $X\trans X / \sigma^2$:
độ lệch chuẩn nhiễu tăng gấp đôi thì mỗi quan sát chỉ mang một phần tư lượng thông tin. Độ lệch
chuẩn hậu nghiệm của hệ số góc với một biến giải thích là xấp xỉ
%
\begin{equation}
  \label{eq:slope-sd}
  \text{sd}(\beta_1 \mid y) \approx
  \Bigl(\frac{n\, s_x^2}{\sigma^2} + \frac{1}{\tau^2}\Bigr)^{-1/2}
  \approx \frac{\sigma}{\sqrt{n}\, s_x},
\end{equation}
%
với $s_x$ là độ lệch chuẩn mẫu của $x$. Muốn giữ nguyên độ chính xác khi $\sigma$
tăng gấp đôi thì cần gấp bốn lần dữ liệu (hình~\ref{fig:noise-spread}, trái).
Khi $\sigma^2$ chưa biết, mức nhiễu cũng được ước lượng từ mẫu qua $b_n$ trong
\eqref{eq:nig} hay qua bước lấy mẫu $\sigma^2$ trong \eqref{eq:gibbs}: tổng bình
phương phần dư lớn thì hậu nghiệm của $\sigma^2$ dịch về phía các giá trị lớn hơn, kéo theo hậu nghiệm
của $\beta$ rộng ra.

\resultfig[0.95]{fig_noise_spread}{Độ lệch chuẩn hậu nghiệm của hệ số góc theo cỡ
  mẫu, với ba mức nhiễu $\sigma$ (trái) và ba mức mức độ phân tán của biến giải thích
  $s_x$ (phải); hai trục theo thang logarit.}{fig:noise-spread}

\section{Mức độ phân tán và phạm vi của biến giải thích}
\label{sec:dep-spread}

Công thức \eqref{eq:slope-sd} cho thấy mức độ phân tán của $x$ cũng ảnh hưởng
đến độ chính xác của ước lượng: $s_x$ gấp bốn lần thì độ lệch chuẩn hậu nghiệm giảm bốn lần, tương đương
gấp 16 lần dữ liệu (hình~\ref{fig:noise-spread}, phải). Dữ liệu tập trung trong
một khoảng hẹp của $x$ cho rất ít thông tin về hệ số góc, dù $n$ lớn. Điều này
ứng với $X\trans X$ có các trị riêng nhỏ.

Phạm vi mà mẫu bao phủ quyết định chất lượng dự đoán. Theo \eqref{eq:predictive},
phần bất định do tham số $x^{\ast\top} V_n x^\ast$ tăng theo bình phương khoảng
cách từ $x^\ast$ tới tâm dữ liệu. Ngoài vùng có dữ liệu, khoảng tin cậy hậu nghiệm của trung
bình mở rộng nhanh và chiếm phần lớn khoảng dự đoán (hình~\ref{fig:predictive}).
Khoảng dự đoán có thể mở rộng khi ngoại suy, phản ánh bất định về tham số
trong mô hình đã chọn. Tuy nhiên, khoảng này không phản ánh đầy đủ bất định về
việc quan hệ tuyến tính có còn phù hợp ngoài phạm vi dữ liệu hay không.

\resultfig[0.8]{fig_predictive}{Khoảng tin cậy hậu nghiệm 95\% của trung bình và khoảng dự
  đoán 95\%, với 20 quan sát có $x$ trong $[0, 1]$; ngoài vùng có dữ liệu, bất
  định về tham số chiếm ưu thế.}{fig:predictive}

\section{Tương quan giữa các biến giải thích}
\label{sec:dep-collinear}

Khi các biến giải thích tương quan mạnh, $X\trans X$ gần suy biến: có những tổ
hợp tuyến tính của hệ số mà dữ liệu gần như không phân biệt được. Hậu nghiệm vì
thế dài ra theo hướng đó, các hệ số tương quan âm với nhau trong hậu nghiệm, và
phương sai biên của từng hệ số tăng lên, dù tổng của chúng vẫn được xác định tốt.
Trong mô phỏng ở hình~\ref{fig:collinearity} với $n = 50$, khi tương quan giữa
$x_1$ và $x_2$ tăng từ 0 lên \num{0.99}, độ lệch chuẩn hậu nghiệm của $\beta_1$
tăng từ \num{0.14} lên \num{0.90} và tương quan hậu nghiệm giữa $\beta_1$ và
$\beta_2$ từ \num{0.06} xuống $-0{,}99$.

\resultfig[0.95]{fig_collinearity}{Hậu nghiệm đồng thời của $(\beta_1, \beta_2)$
  với ba mức tương quan giữa hai biến giải thích; elip là vùng 95\%, dấu chéo là
  giá trị thật.}{fig:collinearity}

Đặc điểm này tác động tới ba bước. Ở bước 3, phân phối tiên nghiệm thêm $V_0^{-1}$ vào $X\trans X$
nên làm ma trận khả nghịch và ổn định hệ số, đúng cơ chế của hồi quy Ridge. Ở
bước 4, hình dạng kéo dài và hẹp của phân phối hậu nghiệm có thể làm giảm hiệu quả
lấy mẫu của Gibbs và Metropolis--Hastings. NUTS có thể hạn chế khó khăn này
nhờ cơ chế điều chỉnh ma trận khối lượng, và có thể giảm
khó khăn tính toán bằng chuẩn hoá hoặc phân rã QR của $X$. Ở bước 7, hệ số riêng lẻ khó diễn
giải, trong khi dự đoán trong phạm vi dữ liệu vẫn có thể ổn định.

\section{Thang đo của biến}
\label{sec:dep-scale}

OLS bất biến với thang đo: đổi đơn vị của $x$ từ kg sang tấn chỉ nhân hệ số với
\num{1000}. Nếu giữ nguyên tiên nghiệm khi đổi đơn vị, kết quả Bayes có thể thay đổi, vì phân phối tiên nghiệm $\beta_j \sim \N(0, \tau^2)$ mang ý
nghĩa khác hẳn khi đơn vị thay đổi: một phân phối tiên nghiệm yếu thông tin theo kg trở thành
tiên nghiệm tập trung hơn nhiều theo tấn. Thang đo còn ảnh hưởng tới bước 4, vì các tham số có độ
lớn rất khác nhau làm các chiều của phân phối hậu nghiệm có thang đo khác nhau,
khiến thuật toán lấy mẫu có thể phải giảm kích thước bước. Chuẩn hoá các
biến giải thích ở bước 1 giúp hạn chế hai vấn đề này: mọi hệ số cùng một thang, nên
có thể dùng chung một phân phối tiên nghiệm và thuật toán lấy mẫu có thể hội tụ nhanh hơn.

\section{Ngoại lai và đuôi dày}
\label{sec:dep-outlier}

Với hàm hợp lý chuẩn, mỗi quan sát đóng góp $-(y_i - x_i\trans\beta)^2 / 2\sigma^2$
vào log-hậu nghiệm; độ lớn của mức phạt tăng theo bình phương phần dư. Một điểm ngoại lai duy nhất
có thể kéo cả đường hồi quy về phía nó, trong khi hậu nghiệm vẫn hẹp. Khi đó, kết quả có thể sai lệch
với mức bất định bị đánh giá thấp. Trong hình~\ref{fig:outlier}, một lỗi lớn ở điểm có $x$ lớn nhất kéo
trung bình hậu nghiệm của hệ số góc từ 2 xuống \num{1.22}. Thay hàm hợp lý chuẩn
bằng Student-t với $\nu = 3$, mức phạt chỉ tăng theo logarit của phần dư khi phần
dư lớn, nên điểm ngoại lai bị giảm trọng số tự động và trung bình hậu nghiệm trở
về \num{1.99}.

\resultfig[0.95]{fig_outlier}{Một điểm ngoại lai: đường hồi quy theo trung bình
  hậu nghiệm (trái) và hậu nghiệm của hệ số góc (phải), với hàm hợp lý chuẩn và
  Student-t. Hậu nghiệm tính trên lưới $(\beta_0, \beta_1)$, $\sigma = 0{,}6$
  cố định.}{fig:outlier}

Ngoại lai tác động tới bước 2 (chọn hàm hợp lý), bước 4 (với Student-t không còn
dạng đóng, phải dùng MCMC) và bước 6, nơi Pareto $k$ lớn của PSIS-LOO và các
thống kê cực trị trong kiểm tra dự đoán hậu nghiệm chỉ ra những quan sát mô hình
không giải thích được.

\section{Dạng quan hệ và phân phối sai số}
\label{sec:dep-shape}

Hàm hợp lý \eqref{eq:model} giả định quan hệ tuyến tính, phương sai không đổi và
sai số độc lập. Mẫu học vi phạm các giả định này thì hậu nghiệm vẫn tính được,
nhưng là hậu nghiệm của một mô hình sai. Các vi phạm thường gặp và cách xử lý:

\begin{itemize}
  \item \textbf{Quan hệ cong hoặc biến phản hồi lệch phải}: biến đổi ở bước 1, như
    lấy logarit, hoặc thêm số hạng bậc hai.
  \item \textbf{Phương sai thay đổi} theo nhóm hay theo mức của $x$: cho $\sigma$
    phụ thuộc nhóm ở bước 2.
  \item \textbf{Sai số phụ thuộc}, như dữ liệu chuỗi thời gian hay dữ liệu lặp lại
    trên cùng một đối tượng: khi đó $n$ quan sát mang ít thông tin hơn $n$ quan
    sát độc lập, và hậu nghiệm từ \eqref{eq:model} hẹp hơn thực tế. Cần thêm cấu
    trúc tự tương quan hoặc hiệu ứng ngẫu nhiên ở bước 2.
  \item \textbf{Dữ liệu thiếu}: có thể loại bỏ các quan sát có giá trị thiếu khi tỉ lệ thiếu
    nhỏ và cơ chế thiếu dữ liệu phù hợp. Một lựa chọn khác là xây dựng mô hình
    cho các giá trị thiếu và ước lượng đồng thời với các tham số của mô hình.
\end{itemize}

\section{Tổng hợp}

Bảng~\ref{tab:summary} tóm tắt mối liên hệ giữa các đặc điểm của mẫu và các đại
lượng trong mô hình, các bước chịu ảnh hưởng và cách xử lý tương ứng.

\begin{table}[tp]
  \centering
  \small
  \caption{Đặc điểm của mẫu học và nơi chúng tác động trong quá trình ước
    lượng.}
  \label{tab:summary}
  \begin{tabular}{>{\raggedright\arraybackslash}p{3.0cm}
                  >{\raggedright\arraybackslash}p{4.4cm}
                  >{\raggedright\arraybackslash}p{2.2cm}
                  >{\raggedright\arraybackslash}p{6.6cm}}
    \toprule
    Đặc điểm & Đại lượng liên quan & Bước & Biểu hiện và cách xử lý \\
    \midrule
    Cỡ mẫu $n$ & $X\trans X$ tỉ lệ với $n$; trọng số $w$ của dữ liệu
      & 3, 4 & Độ lệch chuẩn hậu nghiệm giảm theo $1/\sqrt{n}$; mẫu nhỏ thì phân phối tiên nghiệm quan trọng, cần phân tích độ nhạy \\
    Mức nhiễu $\sigma$ & $X\trans X/\sigma^2$; tổng bình phương phần dư trong $b_n$
      & 4 & Độ lệch chuẩn nhiễu tăng gấp đôi cần cỡ mẫu gấp bốn lần; $\sigma^2$ được ước lượng đồng thời với $\beta$ \\
    Mức độ phân tán $s_x$ & Trị riêng của $X\trans X$
      & 1, 4 & Độ lệch chuẩn hệ số tỉ lệ với $1/s_x$; dữ liệu tập trung cho ít thông tin \\
    Phạm vi của $x$ & $x^{\ast\top} V_n x^\ast$ trong dự đoán
      & 7 & Khoảng dự đoán rộng nhanh khi ngoại suy \\
    Đa cộng tuyến & $X\trans X$ gần suy biến
      & 3, 4, 7 & Hậu nghiệm có dạng kéo dài, các hệ số tương quan; điều chuẩn bằng tiên nghiệm, chuẩn hoá, QR \\
    Thang đo & Ý nghĩa của phân phối tiên nghiệm; hình dạng hậu nghiệm
      & 1, 3, 4 & Chuẩn hoá trước khi đặt phân phối tiên nghiệm \\
    Ngoại lai & Phần dư bình phương trong hàm hợp lý chuẩn
      & 2, 4, 6 & Kéo lệch hệ số; dùng Student-t, kiểm tra Pareto $k$ \\
    Dạng quan hệ, sai số & Giả định của \eqref{eq:model}
      & 1, 2, 6 & Biến đổi biến, mô hình phương sai, tự tương quan; kiểm tra bằng PPC \\
    \bottomrule
  \end{tabular}
\end{table}

% ===========================================================================
\chapter{Minh hoạ trên dữ liệu thực tế}
\label{sec:example}

Bài cuối kỳ của nhóm áp dụng đúng quy trình trên cho 93 mẫu xe mới tại Mỹ năm
1993, với biến phản hồi là hiệu suất nhiên liệu trong thành phố (City, MPG). Một
số kết quả cho thấy các cơ chế ở chương~\ref{sec:dependence} trên dữ liệu thực tế.

\paragraph{Biến đổi và chuẩn hoá (bước 1).} Quan hệ giữa City và công suất có
dạng cong. Lấy logarit City làm logarit hàm hợp lý trên thang City tăng từ $-264{,}0$
lên $-246{,}0$, và phần dư không còn dạng cong. Các biến giải thích được chuẩn hoá
để dùng chung phân phối tiên nghiệm $\beta_j \sim \N(0, 2{,}5^2)$.

\paragraph{Cỡ mẫu và phân phối tiên nghiệm (bước 3).} Với $n = 93$, phân phối tiên nghiệm phẳng và phân phối tiên nghiệm yếu
thông tin cho hậu nghiệm gần như trùng nhau, trung bình lệch nhau tối đa
\num{0.04} độ lệch chuẩn hậu nghiệm. Tiên nghiệm tập trung $\N(0, 0{,}5^2)$ kéo các hệ số về
0 tối đa \num{0.42} độ lệch chuẩn, rõ nhất ở số hạng tương tác vốn được ít dữ
liệu hỗ trợ nhất, đúng như \eqref{eq:shrinkage} dự báo cho hệ số được ước lượng kém chính xác nếu chỉ dựa vào dữ liệu.

\paragraph{Đối chiếu lời giải dạng đóng với MCMC (bước 4, 5).} Hậu nghiệm dạng đóng NIG của hệ
số công suất có trung bình $-0{,}7101$, còn NUTS với cùng phân phối tiên nghiệm cho $-0{,}7099$.
Thuật toán Metropolis--Hastings do nhóm tự triển khai cho cùng kết quả nhưng chỉ đạt cỡ mẫu hiệu dụng
khoảng \num{1700} từ \num{20000} mẫu, so với khoảng \num{12000} từ \num{8000}
mẫu của NUTS.

\paragraph{Tương quan giữa biến giải thích.} Trọng lượng tương quan \num{0.74}
với công suất. Khi thêm trọng lượng vào mô hình, hệ số chuẩn hoá của công suất
thay đổi từ khoảng $-0{,}71$ thành $-0{,}15$. Kết quả này cho thấy mối liên hệ
giữa công suất và hiệu suất nhiên liệu phụ thuộc đáng kể vào việc có đưa trọng
lượng vào mô hình hay không; không thể từ đó kết luận về quan hệ nhân quả.

\paragraph{Ngoại lai (bước 2, 6).} Ba xe cỡ nhỏ có hiệu suất nhiên liệu cao so với phần còn lại của mẫu.
Kiểm tra dự đoán hậu nghiệm cho thấy mô hình sai số chuẩn chưa tái tạo tốt giá trị lớn nhất.
Chuyển sang hàm hợp lý Student-t, độ lệch chuẩn nhiễu hậu nghiệm giảm từ
\num{0.47} xuống \num{0.38}, Pareto $k$ lớn nhất giảm từ \num{0.54} xuống
\num{0.46}, và mô hình này có ELPD ước lượng cao nhất theo PSIS-LOO.

\paragraph{Phạm vi dữ liệu (bước 7).} Mức công suất 300 hp cần dự đoán trùng với
giá trị lớn nhất trong mẫu. Khoảng dự đoán 95\% cho một xe Mỹ 300 hp là 13{,}4 đến
20{,}6 MPG, chứa các giá trị quan sát của ba xe có công suất gần 300 hp; mô hình chỉ dùng công suất ngoại suy
theo đường thẳng và dự đoán thấp hơn thực tế.

% ===========================================================================
\chapter{Kết luận}
\label{sec:conclusion}

Ước lượng mô hình hồi quy tuyến tính Bayes gồm các bước chuẩn bị dữ liệu,
đặc tả hàm hợp lý, chọn tiên nghiệm, tính hậu nghiệm, chẩn đoán hội tụ, kiểm tra
mô hình, suy luận và dự đoán. Trong mô hình sai số chuẩn đang xét, với tiên
nghiệm cố định, dữ liệu ảnh hưởng đến hậu nghiệm thông qua $n$, $X\trans X$,
$X\trans y$ và $y\trans y$.

Cỡ mẫu ảnh hưởng đến vai trò tương đối của tiên nghiệm. Mức nhiễu, mức độ phân
tán và tương quan của biến giải thích ảnh hưởng đến độ chính xác của ước lượng.
Phạm vi dữ liệu giới hạn khả năng đánh giá dự đoán khi ngoại suy, còn thang đo
ảnh hưởng đến cách lựa chọn tiên nghiệm và hiệu quả lấy mẫu. Ngoại lai và vi
phạm giả định có thể dẫn đến kết quả sai lệch với mức bất định bị đánh giá thấp.

Các đặc điểm này cần được xem xét xuyên suốt quy trình: chuẩn hoá và biến đổi
dữ liệu ở bước 1, lựa chọn hàm hợp lý ở bước 2, đặc tả tiên nghiệm và phân tích
độ nhạy ở bước 3, lựa chọn thuật toán ở bước 4, và đánh giá kết quả qua các bước
chẩn đoán, kiểm tra mô hình. Các mô phỏng và ví dụ trên dữ liệu xe minh hoạ mối
liên hệ giữa đặc điểm của mẫu với những lựa chọn này.

% ===========================================================================
\appendix

\chapter{Dẫn xuất hậu nghiệm khi $\sigma^2$ đã biết}
\label{sec:derivation}

Với phân phối tiên nghiệm $\beta \sim \N(m_0, V_0)$ và hàm hợp lý \eqref{eq:model}, bỏ các số hạng
không chứa $\beta$,
%
\begin{equation}
  \label{eq:derive1}
  \log p(\beta \mid y)
  = -\frac{1}{2\sigma^2} \norm{y - X\beta}_2^2
    - \frac{1}{2} (\beta - m_0)\trans V_0^{-1} (\beta - m_0) + \text{const}.
\end{equation}
%
Khai triển và gom các số hạng bậc hai và bậc nhất theo $\beta$,
%
\begin{equation}
  \label{eq:derive2}
  \log p(\beta \mid y)
  = -\frac{1}{2} \beta\trans \Bigl( V_0^{-1} + \tfrac{1}{\sigma^2} X\trans X \Bigr) \beta
    + \beta\trans \Bigl( V_0^{-1} m_0 + \tfrac{1}{\sigma^2} X\trans y \Bigr)
    + \text{const}.
\end{equation}
%
Một phân phối chuẩn $\N(m_n, V_n)$ có log-mật độ
$-\tfrac{1}{2}\beta\trans V_n^{-1}\beta + \beta\trans V_n^{-1} m_n + \text{const}$.
Đồng nhất hệ số của hai biểu thức cho
$V_n^{-1} = V_0^{-1} + X\trans X/\sigma^2$ và
$V_n^{-1} m_n = V_0^{-1} m_0 + X\trans y/\sigma^2$, tức công thức
\eqref{eq:known-sigma}. Vì \eqref{eq:derive2} là dạng toàn phương xác định âm
theo $\beta$, hậu nghiệm đúng là phân phối chuẩn, không chỉ xấp xỉ.

Phân phối điều kiện của $\sigma^2$ trong \eqref{eq:gibbs} có được theo cách tương
tự: với $\beta$ cố định, tích của hàm hợp lý và phân phối tiên nghiệm $\text{IG}(a_0, b_0)$ tỉ lệ
với $(\sigma^2)^{-(a_0 + n/2) - 1} \exp\bigl(-(b_0 + \tfrac{1}{2}\norm{y - X\beta}_2^2)/\sigma^2\bigr)$,
chính là mật độ của $\text{IG}\bigl(a_0 + \tfrac{n}{2},\; b_0 + \tfrac{1}{2}\norm{y - X\beta}_2^2\bigr)$.

\newpage

\printbibliography[heading=bibintoc,title={Tài liệu tham khảo}]

\end{document}
